% !TEX root = master.tex
% !TEX program = lualatex

%%%%%%%%%%%%%%%%%%%%%%%%%%%%%%%%%%%%%%%%%%%%%%%%%%%%%%
\section{Network and File Syscalls}

\subsection{Sockets}

A \textbf{socket} is an abstraction that allows processes to communicate with remote processes.

To the program, a socket is simply an integer.

%%%%%%%%%%%%%%%%%%%%%%%%%%%%%%%%%%%%%%%%%%%%%%%%%%%%%%
\subsection{Creating Communication}

Preparation steps:

\begin{enumerate}
\item socket()
\item bind()
\end{enumerate}

\subsection{Connectionless Communication}

Syscalls:

\begin{itemize}
\item sendto()
\item recvfrom()
\end{itemize}

Characteristics:

\begin{itemize}
\item same socket used for multiple remote processes
\item used by protocols like DNS
\end{itemize}

%%%%%%%%%%%%%%%%%%%%%%%%%%%%%%%%%%%%%%%%%%%%%%%%%%%%%%
\subsection{Connection-Oriented Communication}

Used by web servers.

Server steps:

\begin{enumerate}
\item socket()
\item bind()
\item listen()
\item accept()
\end{enumerate}

Client steps:

\begin{enumerate}
\item socket()
\item connect()
\end{enumerate}

After connection:

\begin{itemize}
\item send()
\item recv()
\end{itemize}

%%%%%%%%%%%%%%%%%%%%%%%%%%%%%%%%%%%%%%%%%%%%%%%%%%%%%%
\subsection{Closing Communication}

When communication ends:

\begin{verbatim}
close(socket)
\end{verbatim}

The socket is destroyed and the port is released.

%%%%%%%%%%%%%%%%%%%%%%%%%%%%%%%%%%%%%%%%%%%%%%%%%%%%%%

